%======================================================================%
\section{Standard Model Descent}
\label{sec:sm}
%======================================================================%

Matter and electroweak symmetry breaking are \textbf{postulated} (P2--P4).
Anomaly cancellation is a \textbf{theorem} once the $\mathbf{27}$ content is
fixed.

%----------------------------------------------------------------------%
\subsection{Postulates P2--P4}
\label{sec:sm:postulates}
%----------------------------------------------------------------------%

\begin{postulate}[Chiral matter, P2]\label{post:sm:P2}
Each family furnishes one chiral $\mathbf{16}$ of $\Spin(10)$ with charges from
\eqref{eq:foundations:27branch}. No mirror $\overline{\mathbf{16}}$ partners.
\end{postulate}

\begin{postulate}[Three families, P3]\label{post:sm:P3}
Horizontal $SU(3)_F$ acts on $\jalg\otimes\C^3$; the three factors are three
chiral $\mathbf{16}$s. Three generations are \emph{not} derived from
Axioms~I--III.
\end{postulate}

\begin{postulate}[GUT--Higgs and Yukawa, P4]\label{post:sm:P4}
Breaking proceeds via $\mathbf{45}_H$ or $\mathbf{54}_H$
($\Spin(10)\to SU(5)\times U(1)_X$), $\mathbf{126}_H$ (breaks $B-L$), and
$\mathbf{10}_H$ (electroweak). Yukawas $y_{ij}\,\mathbf{16}_i\mathbf{16}_j H$
generate fermion masses; type-I seesaw from $\mathbf{126}_H$ gives
$m_\nu\sim m_D^2/M_R$.
\end{postulate}

\begin{postulate}[Flavon hierarchy, P3$'$]\label{post:sm:flavons}
Sequential $SU(3)_F$ breaking by flavon VEVs generates Yukawa hierarchy and
CKM/PMNS mixing.
\end{postulate}

%----------------------------------------------------------------------%
\subsection{Descent chain}
\label{sec:sm:descent}
%----------------------------------------------------------------------%

\begin{theorem}[SM gauge structure under P4]\label{thm:sm:descent}
Assume P4. The breaking chain
\[
  \E{6}\xrightarrow{V}\Spin(10)\times U(1)
  \xrightarrow{\mathbf{45}_H}\mathrm{SU}(5)\times U(1)_X
  \xrightarrow{\mathbf{126}_H}\mathrm{SU}(3)_c\times\mathrm{SU(2)}_L\times U(1)_Y
  \xrightarrow{\mathbf{10}_H}\mathrm{SU}(3)_c\times U(1)_{\mathrm{em}}
\]
yields the SM gauge group with GUT normalization $\alpha_1=(5/3)\alpha_Y$.
\end{theorem}

Per family $\mathbf{16}=(q,u^c,d^c,\ell,e^c,\nu^c)$; $\mathbf{10}_H$ supplies
$v=246\,\mathrm{GeV}$.

\begin{remark}[Yukawa source term]
A bare coupling $\bar\psi\,\orderparam^a T_a\,\psi$ with $\orderparam$ in
$\{\mathbf{16},\overline{\mathbf{16}}\}$ is \emph{not} $H$-invariant and does not
generate masses at the $\orderparam_\star$ vacuum. Masses require P4 Higgs
multiplets~\cite{GeorgiGlashow1974}.
\end{remark}

%----------------------------------------------------------------------%
\subsection{Scalar spectrum}
\label{sec:sm:masses}
%----------------------------------------------------------------------%

\begin{theorem}[Tree-level Goldstones]\label{thm:sm:tree}
At $\orderparam_\star\in\mathcal{M}$, all $32$ coset directions are massless at
tree level (Corollary~\ref{cor:dynamics:goldstone}).
\end{theorem}

\begin{theorem}[Radiative degeneracy]\label{thm:sm:radiative}
The Coleman--Weinberg potential~\cite{ColemanWeinberg1973} is $H$-invariant.
Since $\mathfrak{m}$ is a single real-irreducible $H$-module
(Theorem~\ref{thm:spacetime:irrep}), Schur gives one eigenvalue on all $28$
physical massive scalars; canonical mass $m^2=\mu^2$.
\end{theorem}

\begin{remark}[Heterotic corroboration]\label{rem:sm:heterotic}
The heterotic $\mathbb{Z}_5$ quintic gives $h^1(\hat Q,V_{16})=3$~\cite{Witten1985}.
\textbf{Conjecture}, not a derivation of P3.
\end{remark}